% math0 manuscript. Build with: latexmk -pdf paper.tex
\documentclass[11pt]{amsart}
\usepackage[margin=1.1in]{geometry}
\usepackage{amsmath,amssymb,amsthm,mathtools}
\usepackage{url}
\usepackage{hyperref}
\usepackage{xcolor}
\usepackage{listings}
\usepackage{newunicodechar}

\lstdefinelanguage{lean}{
  keywords={theorem,lemma,def,abbrev,example,fun,by,have,show,from,let,in,if,then,else,match,with,
    namespace,end,open,import,section,variable,instance,structure,where,calc,obtain,intro,exact,
    refine,rcases,rintro,simp,omega,decide,norm_num,linarith,nlinarith,positivity,induction,cases,
    private,set_option,rw,at,push_cast,exact_mod_cast},
  sensitive=true, comment=[l]{--}, morecomment=[s]{/-}{-/}, string=[b]"}
\lstset{language=lean, basicstyle=\ttfamily\footnotesize, columns=fullflexible, keepspaces=true, breaklines=true,
  keywordstyle=\color{blue!60!black}\bfseries, commentstyle=\color{gray}, frame=single,
  extendedchars=true, inputencoding=utf8,
  literate={∑}{{$\sum$}}1 {∈}{{$\in$}}1 {ℕ}{{$\mathbb{N}$}}1 {ℚ}{{$\mathbb{Q}$}}1
    {∀}{{$\forall$}}1 {≥}{{$\ge$}}1 {…}{{$\ldots$}}1 {¬}{{$\neg$}}1 {⊢}{{$\vdash$}}1}
% Unicode commonly found in Lean sources.
\newunicodechar{ℕ}{\ensuremath{\mathbb{N}}}
\newunicodechar{ℤ}{\ensuremath{\mathbb{Z}}}
\newunicodechar{ℚ}{\ensuremath{\mathbb{Q}}}
\newunicodechar{ℝ}{\ensuremath{\mathbb{R}}}
\newunicodechar{ℂ}{\ensuremath{\mathbb{C}}}
\newunicodechar{∀}{\ensuremath{\forall}}
\newunicodechar{∃}{\ensuremath{\exists}}
\newunicodechar{→}{\ensuremath{\to}}
\newunicodechar{←}{\ensuremath{\leftarrow}}
\newunicodechar{↔}{\ensuremath{\leftrightarrow}}
\newunicodechar{↦}{\ensuremath{\mapsto}}
\newunicodechar{¬}{\ensuremath{\neg}}
\newunicodechar{∧}{\ensuremath{\wedge}}
\newunicodechar{∨}{\ensuremath{\vee}}
\newunicodechar{≤}{\ensuremath{\le}}
\newunicodechar{≥}{\ensuremath{\ge}}
\newunicodechar{≠}{\ensuremath{\ne}}
\newunicodechar{∈}{\ensuremath{\in}}
\newunicodechar{∉}{\ensuremath{\notin}}
\newunicodechar{⊆}{\ensuremath{\subseteq}}
\newunicodechar{∩}{\ensuremath{\cap}}
\newunicodechar{∪}{\ensuremath{\cup}}
\newunicodechar{∑}{\ensuremath{\sum}}
\newunicodechar{∏}{\ensuremath{\prod}}
\newunicodechar{⟨}{\ensuremath{\langle}}
\newunicodechar{⟩}{\ensuremath{\rangle}}
\newunicodechar{·}{\ensuremath{\cdot}}
\newunicodechar{∘}{\ensuremath{\circ}}
\newunicodechar{⁻}{\ensuremath{^{-}}}
\newunicodechar{¹}{\ensuremath{^{1}}}
\newunicodechar{²}{\ensuremath{^{2}}}
\newunicodechar{₀}{\ensuremath{_{0}}}
\newunicodechar{₁}{\ensuremath{_{1}}}
\newunicodechar{₂}{\ensuremath{_{2}}}
\newunicodechar{α}{\ensuremath{\alpha}}
\newunicodechar{β}{\ensuremath{\beta}}
\newunicodechar{ε}{\ensuremath{\varepsilon}}
\newunicodechar{δ}{\ensuremath{\delta}}
\newunicodechar{π}{\ensuremath{\pi}}
\newunicodechar{σ}{\ensuremath{\sigma}}
\newunicodechar{φ}{\ensuremath{\varphi}}
\newunicodechar{∞}{\ensuremath{\infty}}
\newunicodechar{⌊}{\ensuremath{\lfloor}}
\newunicodechar{⌋}{\ensuremath{\rfloor}}
\newunicodechar{⌈}{\ensuremath{\lceil}}
\newunicodechar{⌉}{\ensuremath{\rceil}}
\newunicodechar{∣}{\ensuremath{\mid}}
\newunicodechar{≡}{\ensuremath{\equiv}}
\newunicodechar{×}{\ensuremath{\times}}
\newunicodechar{⧸}{\ensuremath{/}}
\newunicodechar{‖}{\ensuremath{\|}}
\newunicodechar{√}{\ensuremath{\surd}}
\newunicodechar{∅}{\ensuremath{\emptyset}}
\newunicodechar{⦃}{\ensuremath{\{\!|}}
\newunicodechar{⦄}{\ensuremath{|\!\}}}
\newunicodechar{𝓝}{\ensuremath{\mathcal{N}}}
\newunicodechar{≫}{\ensuremath{\gg}}
\newunicodechar{≪}{\ensuremath{\ll}}
\newunicodechar{⊢}{\ensuremath{\vdash}}
\newunicodechar{…}{\ensuremath{\ldots}}

\theoremstyle{plain}
\newtheorem{theorem}{Theorem}
\newtheorem{lemma}[theorem]{Lemma}
\newtheorem{proposition}[theorem]{Proposition}
\newtheorem{corollary}[theorem]{Corollary}
\newtheorem{conjecture}[theorem]{Conjecture}
\theoremstyle{definition}
\newtheorem{definition}[theorem]{Definition}
\newtheorem{remark}[theorem]{Remark}

\emergencystretch=3em
% Lean identifiers in running text (breakable at dots and underscores).
\expandafter\def\expandafter\UrlBreaks\expandafter{\UrlBreaks\do\_\do\.}
\DeclareUrlCommand\decl{\urlstyle{tt}}
\newcommand{\N}{\mathbb{N}}
\newcommand{\Q}{\mathbb{Q}}
\newcommand{\Z}{\mathbb{Z}}
\newcommand{\Zl}{\mathbb{Z}_{(11)}}
\DeclareMathOperator{\num}{num}
\DeclareMathOperator{\den}{den}

\title[Lang's formula for OEIS A294512 fails at $n=164$]{A counterexample at $n=164$ to a conjectured formula for the denominators of partial sums of reciprocal octagonal numbers}
\author{math0}
\address{Produced by Claude (an AI model developed by Anthropic) in the math0 project}
\date{October 18, 2026}

\begin{document}
\begin{abstract}
Let $a(n)$ be the denominator, in lowest terms, of $V(n)=\sum_{k=0}^{n}1/((k+1)(3k+1))$ (OEIS A294512), and let $A250400(m)$ be the numerator of the harmonic mean of the first $m$ octagonal numbers.
In 2017 W.~Lang recorded in the OEIS the conjecture that $a(n)=A250400(n+1)/(n+1)$ for all $n\ge 0$.
We show that the conjecture holds at $n$ exactly when $n+1$ is coprime to the numerator of $V(n)$, and that it fails at $n=164$: there $11$ divides both $165$ and the numerator of $V(164)$, and $11\cdot A250400(165)=165\cdot a(164)$, so $A250400(165)/165=a(164)/11$, which is not an integer.
To our knowledge this is a new result, although an elementary one: the formula is still marked as a conjecture in the OEIS, and no earlier counterexample has been recorded.
The refutation is verified in Lean~4 against a statement that math0 formalized from the OEIS entry, whose definitions are checked by the kernel against the listed OEIS terms; this statement is not part of the formal-conjectures repository.
\end{abstract}
\maketitle

\section{Introduction}

The octagonal numbers are $O_j=j(3j-2)$, $j\ge 0$ (OEIS A000567 \cite{A000567}). In Lang's notation \cite{A294512}, the partial sums of their reciprocals are
\[
  V(n)=V(3,1;n)=\sum_{k=0}^{n}\frac{1}{(k+1)(3k+1)}=\sum_{j=1}^{n+1}\frac{1}{O_j},\qquad n\ge 0,
\]
and they converge to $(\pi/\sqrt3+3\log 3)/4=1.2774\ldots$ \cite{A294512}.
Sums of this kind arise from series $\sum_{k\ge1}c(k)/k$ with periodic coefficients, a topic going back to Euler and treated in Koecher's book \cite[pp.~189--193]{Koe}; see the comments in \cite{A294512}.
The OEIS sequence A294512, contributed by Wolfdieter Lang on November~1, 2017 \cite{A294512}, lists the denominators $a(n)$ of $V(n)$ in lowest terms, and states that the corresponding numerators are given in A250401 \cite{A250401}, i.e.\ that $\num V(n)=A250401(n+1)$; we shall see that this identification also fails at $n=164$ (Corollary~\ref{cor:exact}).

Three years earlier, Colin Barker had contributed the sequences A250400 and A250401 \cite{A250400,A250401}: the numerators and denominators of the harmonic mean $H(m)=m/\sum_{j=1}^{m}O_j^{-1}=m/V(m-1)$ of the first $m$ octagonal numbers, with tables for $m\le 1000$.
Because $H(n+1)=(n+1)/V(n)$, one expects the numerator of $H(n+1)$ to be $(n+1)\,a(n)$.
Lang recorded this in the FORMULA section of A294512 as a conjecture, $a(n)=A250400(n+1)/(n+1)$ for $n\ge 0$, and added a comment to A250401 that assumes it (``assuming that $n+1$ divides $A250400(n+1)$'').
The formula holds for the $23$ listed terms $n\le 22$, and in fact for all $n\le 163$.

\subsection*{Result}
The conjecture holds at $n$ if and only if $\gcd(n+1,\num V(n))=1$ (Lemma~\ref{lem:reduction}), and this fails at $n=164$, where $11$ divides $165$ and $\num V(164)$ (Lemma~\ref{lem:11adic}). The formal proof checks $11\cdot A250400(165)=165\cdot a(164)$ by kernel evaluation and derives a contradiction (Theorem~\ref{thm:main}). Below $5000$ the failures are $n=164$, $2413$, $4520$ (Remark~\ref{rem:more}; not formalized).

\subsection*{Status and prior work}
The conjecture still carries the tag ``[conjecture]'' in the current revision of A294512 (revision~\#13, November~12, 2017), and the comment in A250401 is still stated conditionally.
The math0 referees read the edit histories of A294512, A250400 and A250401 and checked the sibling entries A294513 and A294515; none records a counterexample.
They also found no discussion on the SeqFan mailing list, MathOverflow, Mathematics Stack Exchange or arXiv.
The conjecture is not among the OEIS statements of the formal-conjectures repository \cite{FC}, nor among the $492$ OEIS conjectures, formalized in Lean by Tsoukalas et al., that make up the OEIS Open benchmark \cite{Ad}.
We therefore classify the result as \emph{new}, with two caveats.
First, it is elementary and of modest significance: the counterexample is found by exact rational arithmetic over $165$ terms.
Second, the decisive value was implicitly available since 2014: Barker's table of A250400 \cite{A250400} lists $A250400(165)$, which equals $15\,a(164)$; what was apparently never done is to compare it with the denominator of $V(164)$, since the table of A294512 stops at $n=22$.

\section{Statement}\label{sec:statement}

For $x\in\Q$ we write $x=\num(x)/\den(x)$ in lowest terms with $\den(x)\ge 1$.

\begin{definition}\label{def:main}
For $n\ge 0$ and $m\ge 1$ let
\begin{gather*}
  V(n)=\sum_{k=0}^{n}\frac{1}{(k+1)(3k+1)},\qquad a(n)=\den V(n),\\
  H(m)=\frac{m}{\sum_{j=1}^{m}1/(j(3j-2))},\qquad A250400(m)=\num H(m).
\end{gather*}
\end{definition}

Thus $a$ is A294512 (offset $0$) and $A250400$ is A250400 (offset $1$). The FORMULA section of A294512 (W.~Lang, Nov.~01, 2017) contains, verbatim:
\begin{quote}
\texttt{a(n) = A250400(n+1)/(n+1), n >= 0. [conjecture].}
\end{quote}

\begin{conjecture}[Lang \cite{A294512}]\label{conj}
For every $n\ge 0$, $A250400(n+1)/(n+1)=a(n)$, the quotient being taken in $\Q$.
\end{conjecture}

\begin{theorem}\label{thm:main}
$11\cdot A250400(165)=165\cdot a(164)$. Consequently $A250400(165)/165=a(164)/11\ne a(164)$, and Conjecture~\ref{conj} is false.
\end{theorem}

In addition, $11\nmid a(164)$ (Corollary~\ref{cor:exact}), so $A250400(165)/165$ is not even an integer. Since $a(164)\ge1$, the conjecture also fails at $n=164$ if the quotient is read as integer division: $\lfloor a(164)/11\rfloor<a(164)$.

\subsection*{Lean formalization}
The statement was formalized by math0 in the file \path{lean/Math0/Statements/OEIS/A294512.lean}; it was reviewed for fidelity and frozen (SHA-256 \texttt{20fd25a6\ldots f89848}) before any proof attempt. Its definitions and the statement are, verbatim except that docstrings and three data lemmas are omitted:
\begin{lstlisting}
namespace Math0.Statements.A294512

def V (n : ℕ) : ℚ := ∑ k ∈ Finset.range (n + 1), 1 / (((k : ℚ) + 1) * (3 * (k : ℚ) + 1))

def a (n : ℕ) : ℕ := (V n).den

def octagonal (j : ℕ) : ℕ := j * (3 * j - 2)

def harmonicMean (m : ℕ) : ℚ := (m : ℚ) / ∑ j ∈ Finset.Icc 1 m, 1 / (octagonal j : ℚ)

def A250400 (m : ℕ) : ℕ := (harmonicMean m).num.natAbs

def statement : Prop :=
  ∀ n : ℕ, (A250400 (n + 1) : ℚ) / ((n : ℚ) + 1) = (a n : ℚ)

end Math0.Statements.A294512
\end{lstlisting}
The encoding matches Definition~\ref{def:main}.
\decl{V} is the rational sum $V(n)$, and \decl{a} is its reduced denominator \decl{Rat.den}, which is positive.
\decl{octagonal} is $j(3j-2)$ computed in $\N$; the truncated subtraction $3j-2$ is exact for $j\ge1$, and only $j\in\{1,\dots,m\}$ enters \decl{harmonicMean}.
\decl{A250400} is the reduced numerator \decl{Rat.num} of the harmonic mean, an integer that is positive for $m\ge1$, so \decl{Int.natAbs} returns it unchanged.
The statement quantifies over all $n\in\N$, i.e.\ $n\ge0$, and divides in $\Q$ by $n+1\ne0$, so it asserts that $n+1$ divides $A250400(n+1)$ with quotient $a(n)$.

The file also contains three lemmas, each proved by \texttt{decide +kernel}, that compare the definitions with the OEIS data: \decl{a_data} checks all $23$ listed terms $a(0),\dots,a(22)$ of A294512 (from $1,8,168,420,\dots$ to $973748862277067676800$), \decl{A250400_data} checks all $20$ listed terms $A250400(1),\dots,A250400(20)$ (from $1,16,504,\dots$ to $162446703896568000$), and \decl{octagonal_data} checks $O_0,\dots,O_4=0,1,8,21,40$ against A000567.

\section{Proof}\label{sec:proof}

The formal proof has two steps: a kernel evaluation of the identity $11\cdot A250400(165)=165\cdot a(164)$ (Proposition~\ref{prop:key}), and a short deduction of the negation of the statement from it (Section~\ref{sec:deduction}). Before giving them we explain why the identity holds; Lemmas~\ref{lem:reduction} and~\ref{lem:11adic} are not part of the formal proof, but together they give a second, hand-checkable proof that the conjecture fails at $n=164$.

\subsection{Reduction to a coprimality condition}

For $n\ge0$ write $V(n)=p_n/q_n$ in lowest terms, so $p_n,q_n\ge1$, $\gcd(p_n,q_n)=1$ and $q_n=a(n)$, and put $g_n=\gcd(n+1,p_n)$.

\begin{lemma}\label{lem:reduction}
For every $n\ge0$ we have $A250400(n+1)=(n+1)\,a(n)/g_n$, that is,
\[
  g_n\cdot A250400(n+1)=(n+1)\cdot a(n),\qquad\text{and}\qquad \den H(n+1)=p_n/g_n.
\]
In particular Conjecture~\ref{conj} holds at $n$ if and only if $\gcd(n+1,\num V(n))=1$.
\end{lemma}

\begin{proof}
Substituting $j=k+1$ gives $(k+1)(3k+1)=j(3j-2)$, so $\sum_{j=1}^{n+1}1/O_j=V(n)$ and $H(n+1)=(n+1)/V(n)=(n+1)q_n/p_n$.
Since $\gcd(q_n,p_n)=1$, we have $\gcd((n+1)q_n,p_n)=\gcd(n+1,p_n)=g_n$. So the reduced form of $H(n+1)$ has numerator $(n+1)q_n/g_n$ and denominator $p_n/g_n$, both positive. Hence $A250400(n+1)=(n+1)q_n/g_n$ and $\den H(n+1)=p_n/g_n$, so $A250400(n+1)/(n+1)=a(n)/g_n$, which equals $a(n)$ if and only if $g_n=1$.
\end{proof}

\subsection{The prime $11$ divides the numerator of $V(164)$}

Let $\Zl$ denote the ring of rationals whose reduced denominator is not divisible by $11$. For $x,y\in\Zl$ and $e\ge1$ we write $x\equiv y\pmod{11^e}$ if $x-y\in 11^e\Zl$.

\begin{lemma}\label{lem:11adic}
$V(164)\in 11\Zl$. Equivalently, $11\mid p_{164}$; and then $11\nmid q_{164}=a(164)$.
\end{lemma}

\begin{proof}
From $\frac{3}{3j-2}-\frac{1}{j}=\frac{2}{j(3j-2)}$ we get $2V(164)=\sum_{j=1}^{165}\bigl(\frac{3}{3j-2}-\frac1j\bigr)$.
Since $3\cdot 8-2=22$, we have $11\mid 3j-2$ if and only if $j\equiv 8\pmod{11}$. Let $J_1=\{8+11t:0\le t\le 14\}$ be the set of such $j\le165$, and let $J_2=\{11i:1\le i\le15\}$ be the set of multiples of $11$ up to $165$. For $j=8+11t$ we have $3j-2=11(3t+2)$, and $j=11i$ with $i=t+1$ runs over $J_2$ as $t$ runs over $0,\dots,14$. Hence
\[
  \sum_{j\in J_1}\frac{3}{3j-2}-\sum_{j\in J_2}\frac1j=\frac1{11}\sum_{t=0}^{14}\Bigl(\frac{3}{3t+2}-\frac{1}{t+1}\Bigr)=\frac{W}{11},
  \qquad W=\sum_{t=0}^{14}\frac{1}{(t+1)(3t+2)},
\]
and therefore
\[
  2V(164)=U+\frac{W}{11},\qquad U=\sum_{\substack{j\le165\\ j\notin J_1}}\frac{3}{3j-2}-\sum_{\substack{j\le165\\ j\notin J_2}}\frac1j\in\Zl.
\]

\emph{The term $U$.} As $j$ runs over $1,\dots,165$, it runs $15$ times through a complete residue system modulo $11$, and so does $3j-2$, since $j\mapsto 3j-2$ is a bijection of $\Z/11\Z$. Hence
\[
  U\equiv 45\sum_{r=1}^{10}r^{-1}-15\sum_{r=1}^{10}r^{-1}=30\sum_{r=1}^{10}r^{-1}\pmod{11}.
\]
As $r\mapsto r^{-1}$ permutes the nonzero residues modulo $11$, $\sum_{r=1}^{10}r^{-1}\equiv\sum_{r=1}^{10}r=55\equiv0\pmod{11}$. Thus $U\in 11\Zl$.

\emph{The term $W$.} This is Lang's sum $V(3,2;14)$ from the sibling entry A294513 \cite{A294513}. Its $15$ terms have denominators $(t+1)(3t+2)=2,10,24,44,70,102,140,184,234,290,352,420,494,574,660$. Only the terms with $t=3,10,14$ have a denominator divisible by $11$, and these three sum to an element of $\Zl$:
\[
  \frac1{44}+\frac1{352}+\frac1{660}=\frac1{11}\Bigl(\frac14+\frac1{32}+\frac1{60}\Bigr)=\frac{1}{11}\cdot\frac{143}{480}=\frac{13}{480}.
\]
So $W\in\Zl$. In fact $W$ is divisible by $11^2$ in $\Zl$: exact summation gives
$W=832564679309/1157486934240$ in lowest terms, where $832564679309=11^2\cdot31\cdot59\cdot3762001$ and $1157486934240=2^5\cdot3^2\cdot5\cdot7\cdot13\cdot17\cdot19\cdot23\cdot29\cdot41$. (The OEIS lists this numerator and denominator as $A250328(15)$ and $A294513(14)$ \cite{A250328,A294513}.) Hence $W/11\in 11\Zl$.

Combining, $2V(164)\in11\Zl$, and since $2$ is a unit in $\Zl$, $V(164)\in 11\Zl$. Writing $V(164)=11c/d$ with $c,d\in\Z$ and $11\nmid d$, we get $p_{164}d=11cq_{164}$, so $11\mid p_{164}$; and $11\nmid q_{164}$ because $\gcd(p_{164},q_{164})=1$.
\end{proof}

Lemmas~\ref{lem:reduction} and~\ref{lem:11adic} already refute Conjecture~\ref{conj}: $g_{164}=\gcd(165,p_{164})$ is divisible by $11$, so $A250400(165)/165=a(164)/g_{164}$ differs from $a(164)$. To find $g_{164}$ exactly one needs $p_{164}$ modulo $3$ and $5$.

\subsection{The certificate}

Exact rational arithmetic gives $V(164)=p/q$ with $p=p_{164}$ and $q=q_{164}=a(164)$ equal to the following $157$-digit integers (each printed on two lines):
{\small
\[
\begin{aligned}
p={}&1942293989039716766940869656654554278426371164069471388907436248366358029686834\\
    &113509861331525862192731314220412361222824903811876016385114687768467695418619,\\[2pt]
q={}&1522900934001220110486931277055001102268627529692987106334059289643477906495044\\
    &831840380532515782766923778283606417499970869839538732382536890476213388992000.
\end{aligned}
\]}
Here $p\equiv1\pmod 3$, $p\equiv4\pmod5$, $p\equiv0\pmod{11}$, $p\not\equiv0\pmod{121}$, and $q\equiv5\pmod{11}$.

\begin{proposition}\label{prop:key}
$A250400(165)\cdot 11=165\cdot a(164)$.
\end{proposition}

\begin{proof}
\emph{Formal proof.} This is an equality of two closed natural-number expressions. In Lean it is the private theorem \decl{key}, proved by \texttt{decide +kernel}: the kernel evaluates both sides directly from the definitions of Section~\ref{sec:statement} and confirms that they coincide (Section~\ref{sec:lean}).

\emph{Explanation.} By the residues above, $g_{164}=\gcd(3\cdot5\cdot11,p)=11$, so Lemma~\ref{lem:reduction} gives $11\cdot A250400(165)=165\cdot a(164)$.
\end{proof}

\begin{corollary}\label{cor:exact}
$A250400(165)=15\,a(164)$, $\den H(165)=\num V(164)/11$, and $11\nmid a(164)$. In particular $A250400(165)/165=a(164)/11\notin\Z$.
\end{corollary}

\begin{proof}
The first equality is Proposition~\ref{prop:key} divided by $11$; the second is Lemma~\ref{lem:reduction} with $g_{164}=11$; the third is Lemma~\ref{lem:11adic}.
\end{proof}

The value $A250400(165)=15q$, a $158$-digit number, agrees with Barker's table \cite{A250400} for $m\le1000$, which two math0 referees checked independently against Definition~\ref{def:main} for all $m\le 1000$. By the second equality, $A250401(165)=p/11$ rather than $p$, so the comment of Lang in A250401 that identifies $A250401(n+1)$ with the numerator of $V(n)$ also fails at $n=164$, as its stated assumption does; so does the comment in A294512 that the corresponding numerators are given in A250401.

\subsection{Proof of Theorem~\ref{thm:main}}\label{sec:deduction}
This follows the Lean proof of \decl{result} step by step. Suppose Conjecture~\ref{conj} holds, and take $n=164$: $A250400(165)/165=a(164)$ in $\Q$. Multiplying by $165\ne0$ gives $A250400(165)=165\cdot a(164)$ in $\Q$, and hence in $\N$, because the map $\N\to\Q$ is injective. By Proposition~\ref{prop:key}, $11\cdot165\cdot a(164)=11\cdot A250400(165)=165\cdot a(164)$, so $1650\cdot a(164)=0$ and $a(164)=0$. But $a(164)=\den V(164)\ge1$. This contradiction shows that Conjecture~\ref{conj} is false. The relation $A250400(165)/165=a(164)/11$ follows from Proposition~\ref{prop:key} by dividing by $11\cdot165$, and it differs from $a(164)$ because $a(164)\ne0$.\qed

\begin{remark}\label{rem:more}
The following facts come from exact rational computations in Python and are not formally verified.
For $0\le n<5000$, $\gcd(n+1,\num V(n))>1$ exactly for $n=164$, $2413$ and $4520$, with gcd $11$, $17$ and $137$; so $n=164$ is the least counterexample.
Lang's analogous conjecture in A294513 \cite{A294513}, for $V(3,2;n)=\sum_{k\le n}1/((k+1)(3k+2))$ and pentagonal numbers, holds for $n<2500$ except at $n=1891$, where $11$ divides both $1892$ and $\num V(3,2;1891)$.
\end{remark}

\section{Machine verification}\label{sec:lean}

\subsection*{Files and declarations}
The proof is in the file \path{lean/Math0/Frontier/OEIS/A294512.lean} of the math0 repository, \url{https://github.com/Foundation-0/math0}. It imports the frozen statement file \path{lean/Math0/Statements/OEIS/A294512.lean} (Section~\ref{sec:statement}) and the audit commands in \path{lean/Math0/Audit.lean}. Its contents after the module docstring are, verbatim except that docstrings are omitted:
\begin{lstlisting}
namespace Math0.Frontier.A294512
open Math0.Statements.A294512

set_option maxHeartbeats 4000000 in
private theorem key : A250400 165 * 11 = 165 * a 164 := by
  decide +kernel

theorem result : ¬ statement := by
  intro h
  have h1 := h 164
  have hk := key
  have hpos : 0 < a 164 := (V 164).den_pos
  have h2 : (A250400 165 : ℚ) = 165 * (a 164 : ℚ) := by
    rw [div_eq_iff (by norm_num)] at h1
    push_cast at h1 ⊢
    linarith
  have h3 : A250400 165 = 165 * a 164 := by exact_mod_cast h2
  omega

end Math0.Frontier.A294512

#math0_settles Math0.Frontier.A294512.result Math0.Statements.A294512.statement
\end{lstlisting}
The main declaration \decl{Math0.Frontier.A294512.result} has type $\neg$\,\decl{Math0.Statements.A294512.statement}; its proof is that of Section~\ref{sec:deduction}, with \decl{omega} doing the final arithmetic.

\subsection*{Kernel computation}
The lemma \decl{key} is the certificate of Proposition~\ref{prop:key}. With \texttt{decide +kernel}, the Lean kernel itself reduces the decision procedure for this equality in $\N$ to \decl{true}: it unfolds the definitions of Section~\ref{sec:statement}, adds the $165$ rationals of each of the two sums, normalizing to lowest terms with its built-in arbitrary-precision arithmetic on natural-number literals, and compares the two resulting $159$-digit numbers. No compiled code is trusted; \decl{native_decide}, which would add the axiom \decl{Lean.ofReduceBool}, is not used. The option \texttt{maxHeartbeats 4000000} applies to \decl{key} only. The whole file compiles in about $20$ seconds on our machine, most of it spent loading Mathlib.

\subsection*{Audit}
The file ends with \texttt{\#math0\_settles}, defined in \path{lean/Math0/Audit.lean} for statements that math0 formalizes itself. It fails unless (1) \decl{statement} is a definition of type \texttt{Prop} without universe parameters; (2) \decl{statement} and \decl{result} depend on no axioms other than \decl{propext}, \decl{Classical.choice} and \decl{Quot.sound} (this rejects \decl{sorryAx}, \decl{Lean.ofReduceBool}, \decl{Lean.trustCompiler} and custom axioms); and (3) the type of \decl{result} is syntactically the statement or its negation, after removing metadata annotations, with no unfolding. On this file it reports
\begin{lstlisting}
math0 audit: 'Math0.Statements.A294512.statement' depends on axioms [Classical.choice, Quot.sound, propext]
math0 audit: 'Math0.Frontier.A294512.result' depends on axioms [Classical.choice, Quot.sound, propext]
math0 statement check: 'Math0.Frontier.A294512.result' refutes 'Math0.Statements.A294512.statement'
\end{lstlisting}

\subsection*{Toolchain and reproduction}
We use Lean~4 \cite{Lean4} \texttt{v4.33.1} and Mathlib \cite{Mathlib} \texttt{v4.33.1}, pinned in \path{lean/}; the repository also pins formal-conjectures \cite{FC} at commit \texttt{c77fdd4}, which this result does not use.
To reproduce, run \path{scripts/verify.sh} from the repository root: it fetches the Mathlib build cache, builds all math0 results with \texttt{lake build} (re-running every audit command), checks that the pinned \texttt{formal\_conjectures} and \texttt{mathlib} checkouts are unmodified, and reports the number of passed statement checks. The single file can be checked with \texttt{lake env lean Math0/Frontier/OEIS/A294512.lean} in the directory \path{lean}.

\begin{thebibliography}{99}

\bibitem{Ad}
T.~Adamczewski, \emph{OEIS Open: How many conjectures can language models turn into theorems?}, arXiv:2608.11941, 2026.

\bibitem{FC}
M.~Firsching, P.~Lezeau, S.~Mercuri, M.~Z. Horv\'ath, Y.~Dillies, C.~S\"onne, E.~Wieser, F.~Zhang, T.~Hubert, B.~Ag\"uera y Arcas, and P.~Kohli, \emph{Formal Conjectures: An open and evolving benchmark for verified discovery in mathematics}, arXiv:2605.13171, 2026. Repository: \url{https://github.com/google-deepmind/formal-conjectures}.

\bibitem{Koe}
M.~Koecher, \emph{Klassische elementare Analysis}, Birkh\"auser, Basel, Boston, 1987.

\bibitem{Mathlib}
The mathlib Community, \emph{The Lean mathematical library}, in: Proceedings of the 9th ACM SIGPLAN International Conference on Certified Programs and Proofs (CPP 2020), ACM, 2020, 367--381.

\bibitem{Lean4}
L.~de Moura and S.~Ullrich, \emph{The Lean 4 theorem prover and programming language}, in: Automated Deduction -- CADE 28, Lecture Notes in Comput. Sci. \textbf{12699}, Springer, 2021, 625--635.

\bibitem{A000567}
The OEIS Foundation Inc., \url{https://oeis.org/A000567} (octagonal numbers).

\bibitem{A250328}
The OEIS Foundation Inc., \url{https://oeis.org/A250328}, entry by C.~Barker (2014).

\bibitem{A250400}
The OEIS Foundation Inc., \url{https://oeis.org/A250400}, entry by C.~Barker (2014).

\bibitem{A250401}
The OEIS Foundation Inc., \url{https://oeis.org/A250401}, entry by C.~Barker (2014), comment by W.~Lang (2017).

\bibitem{A294512}
The OEIS Foundation Inc., \url{https://oeis.org/A294512}, entry by W.~Lang (2017).

\bibitem{A294513}
The OEIS Foundation Inc., \url{https://oeis.org/A294513}, entry by W.~Lang (2017).

\end{thebibliography}

\end{document}
